% authority: supporting
% status: draft/control
% task: RT-20260718-036
\section{Fresh smuggling audit of the cycle--boundary selector candidate}

\paragraph{Scope.}
This is one bounded audit of the exact proposal-only object
\(\mathrm{EqSrcCycleBoundarySelectorLaw}_{\rm src}^{\rm cand,v1}\), whose
source hash is
\texttt{64901302bd6f37819ebd30aaaebaf39611d2c71c97d45b6cea6a39df11e0df94}.
It supplies no candidate repair, Refuter stress, ontology adoption,
physical-admissibility claim, general \(\mathrm{EqSrc}\) discharge, or
downstream GR promotion.

\subsection{Candidate dependency tree}

An admitted record is \(S=(X,P)\), where \(X\) is finite, \(P\subset X\) is
nonempty and proper, and both \(|X|\) and \(|P|\) are constrained so that
\(|P|\) and \(|Q|\), \(Q=X\setminus P\), are odd.  Put
\[
 p=\mathbf 1_P,\qquad q=\mathbf 1_Q,\qquad
 A_S=\operatorname{span}_{\mathbb F_2}\{p,q\},\qquad
 B_S=\mathbb F_2q .
\]
The quotient label is
\[
 \pi_S(\alpha p+\beta q)=\alpha ,
\]
and the candidate relation is its kernel pair.  The declared variations are
\[
 J_S(\epsilon)(z)=z+\epsilon q,\qquad \epsilon\in\mathbb F_2 .
\]
Thus the displayed dependency order is
\[
 (X,P)\longrightarrow(p,q)\longrightarrow J_S,\ A_S,\ B_S
 \longrightarrow \pi_S\longrightarrow\mathrm{EqCand}_S .
\]

No target atlas, target metric, benchmark success, empirical detector,
workflow state, registry value, role, validator, handoff, checkpoint, or
generated derivative occurs as a mathematical premise in this displayed
tree.  The candidate is therefore source-factorized as written.  This
statement is syntactic and provenance-sensitive: it does not prove that an
actual supplied mark \(P\) was selected without looking at a desired target.

\subsection{Transparent loading and provenance non-identifiability}

The oriented mark \(P\), the odd-cardinality restriction, the field
\(\mathbb F_2\), and the parity functional are explicit assumptions.  Current
ontology does not derive them.  Their explicit appearance is transparent
assumption loading, not hidden target or process smuggling.

Nevertheless the final record does not identify the provenance of \(P\).
Two histories---one in which \(P\) is obtained from an independently
source-side rule and one in which \(P\) is chosen to produce a desired target
relation---yield the same mathematical tuple \((X,P)\) and the same downstream
objects.  Therefore absence of a displayed target premise is not a provenance
certificate.  The no-target-import guard must remain an external acquisition
obligation until a source-only mark law is supplied.

\subsection{Bare-carrier naturality obstruction}

\paragraph{Lemma (no nontrivial natural mark on bare finite carriers).}
There is no assignment \(X\mapsto P_X\) of a nonempty proper subset to every
finite set that is equivariant under all bijections.

\paragraph{Proof.}
Equivariance under all bijections implies that \(P_X\) is invariant under the
full symmetric group \(\operatorname{Sym}(X)\).  This action is transitive.
If \(P_X\) contains one element, invariance carries it to every element, so
\(P_X=X\); if it contains none, \(P_X=\varnothing\).  Hence only the empty and
full subsets are natural on a bare finite carrier.  Both are excluded by the
candidate domain. \(\square\)

The oriented mark must therefore come from richer source data or a new
primitive.  Mark-preserving naturality is valid only after that extra
decoration is supplied.

\subsection{Exact affine variation census}

Write each element of \(A_S\) as the coordinate vector
\((\alpha,\beta)\in\mathbb F_2^2\) in the basis \((p,q)\).  The kernel-pair
relation is equality of the first coordinate.  The affine group
\(\operatorname{AGL}(2,2)\) has
\[
 |\operatorname{GL}(2,2)|\,|A_S|=6\cdot4=24
\]
elements.

\paragraph{Proposition (relation and label stabilizers).}
Among the 24 affine bijections \(F(z)=Lz+t\):
\[
\begin{array}{c|c}
\text{class} & \text{cardinality}\\ \hline
\text{relation stabilizers} & 8\\
\text{pointwise quotient-label fixers} & 4\\
\text{declared candidate variations }J_S & 2\\
\text{relation-changing affine bijections} & 16 .
\end{array}
\]

\paragraph{Proof.}
An affine bijection preserves the kernel-pair relation exactly when
\(L(B_S)=B_S\).  In \(\operatorname{GL}(2,2)\), the stabilizer of the
one-dimensional subspace \(B_S=\mathbb F_2q\) has size two.  The translation
part may be any of four vectors, so there are \(2\cdot4=8\) relation
stabilizers and \(24-8=16\) relation changers.

Pointwise label preservation additionally requires \(t\in B_S\), leaving
\(2\cdot2=4\) maps.  The candidate declares only the two translations
\(\{z\mapsto z,z\mapsto z+q\}\), so it is a proper subset of both the
pointwise-label fixer and relation-stabilizer classes. \(\square\)

Translation by \(p\) is an exact label--relation separation witness:
\[
 \pi_S(z+p)=\pi_S(z)+1
\]
for every \(z\), while equality of quotient labels, and hence the relation,
is preserved for every ordered pair.  The candidate's declared \(J_S\) is
therefore not the full relation-preserving affine variation class.  Its
physical completeness or privileged status is unproved.

\paragraph{Proposition (three-kernel selector nonuniqueness).}
The four-state space \(A_S\cong\mathbb F_2^2\) has exactly three nonzero
proper subspaces,
\[
 \mathbb F_2p,\qquad \mathbb F_2q,\qquad \mathbb F_2(p+q).
\]
The marked automorphism group fixes \(A_S\) pointwise, so all three lines are
invariant.  Each line \(B\) defines a quotient kernel pair and a matching
translation class \(\{z\mapsto z+b:b\in B\}\) with the same acyclic
definition pattern and marked naturality.  Hence marked symmetry does not
select \(B_S=\mathbb F_2q\).  The parity functional supplies substantive
selector data.

\subsection{Complement-mark family}

\paragraph{Theorem (orientation sensitivity on every admitted even carrier).}
Let \(X\) be finite of even cardinality and let \(P\subset X\) be nonempty,
proper, and odd.  Then \(Q=X\setminus P\) is odd, so both \(S=(X,P)\) and
\(S^\complement=(X,Q)\) are admitted.  They satisfy
\[
 A_S=A_{S^\complement}=\operatorname{span}\{p,q\},
 \qquad B_S=\mathbb F_2q,\qquad
 B_{S^\complement}=\mathbb F_2p .
\]
Since \(p,q\) are nonzero with disjoint support, these one-dimensional
subspaces are distinct.  Hence the two orientations induce different
relations on the same four-state space \(A_S\).

\paragraph{Proof.}
Even minus odd is odd, so both marks meet the parity conditions.  Swapping
the displayed mark swaps \(p\) and \(q\), leaving their span fixed while
swapping which line is selected as \(B\).  Distinct nonzero disjoint-support
vectors cannot span the same line over \(\mathbb F_2\). \(\square\)

The four-point example in the candidate is the smallest nondegenerate member
of this family, not an isolated accident.  It proves that orientation is
substantive rather than notational.  It does not prove general EqSrc,
continuum applicability, empirical equivalence, or exact-GR recovery.

\subsection{Naturality, covariance, and robustness boundaries}

Marked-bijection naturality is mathematically exact in the purpose-built
groupoid whose morphisms satisfy \(f(P)=P'\).  Such morphisms exclude the
mark-changing and mark-forgetting transformations most capable of challenging
the selector.  The theorem is covariance of the decorated presentation, not
physical covariance under an independently derived gauge category.

Defining \(J_S\) before \(B_S\) removes literal definitional circularity.
It does not create justificatory independence: both the variation generator
\(q\) and the boundary line \(B_S=\mathbb F_2q\) are engineered from the same
selector-bearing root \(P\).  Robustness is exact because every declared move
lies in the quotient kernel.  That conditional theorem does not establish
that the declared moves are physically admissible, exhaustive, or selected
independently of the relation they preserve.

Every admitted record also yields the same two-dimensional \(A_S\) and
one-dimensional quotient.  This is a uniform one-bit algebraic model.  It is
not evidence that carrier size and structure are physically irrelevant, nor
does the terminology ``cycle'' and ``boundary'' supply missing topological or
homological semantics.

\subsection{Audit classification}

The exact audit result is
\[
\texttt{conditional\_source\_purity\_pass\_with\_transparent\_selector\_loading}
\]
with the mandatory limits
\[
\texttt{mark\_provenance\_underived},\quad
\texttt{variation\_scope\_underived},\quad
\texttt{physical\_covariance\_unproved}.
\]
No forbidden source appears in the displayed definition, and the candidate's
formal theorems survive this audit.  The oriented mark is a transparent,
noneliminable new primitive rather than a current-ontology derivation.

The source-extension classification remains
\texttt{new\_ontology\_primitive\_candidate}; status remains
\texttt{proposal-only}, \texttt{source-extension data}, and
\texttt{blocked\_adoption\_open\_continuation}.  General \(\mathrm{EqSrc}\),
the exact-GR benchmark, and the Distance-to-GR and metric-use ledgers are
unchanged.

\subsection{Freeze and exact next route}

The exact v3 route remains locally frozen.  Repetition of this construction
or this same audit without new mathematical payload is frozen after this
packet.  The marked-selector route is not frozen for exactly one materially
distinct adversarial stress.

The next lawful route is one bounded
\texttt{refuter@0.2.0} \texttt{ontology-law-research-packet} against the exact
candidate hash above.  It should stress mark eliminability, complement and
forgetful transformations, the necessity of odd parity and \(\mathbb F_2\),
independence of physical variations from the selected kernel, the exact
\(8/4/2/16\) affine census, and the limits of the uniform one-bit quotient.
It may not repair, adopt, or promote the candidate.

\paragraph{Forbidden conclusions.}
This audit does not authorize a canonical ontology edit, source-law adoption,
physical admissibility or covariance, general EqSrc, downstream GR objects,
ledger change, benchmark promotion, Gate Chair verdict, completed derivation,
future source-extension impossibility, or global theory rejection.
